% GENERATED FILE. Edit canonical lessons, then run npm run generate:books.
% canonical-source-hash: fnv1a64:7f3abfcfa6993c8a
% canonical-lessons: HI-C489-gholna, HI-C489-catana, HI-C489-khurak, HI-C489-marham, HI-C489-karha

\chapter{Medicines}
\label{ch:hi-medicines}
\hlchaptermodality{489}

\begin{chapteropening}
\textbf{By the end of this chapter:} \emph{I can say how to take or use a medicine.}

Complete the last lesson of chapter 489: I can say how to take or use a medicine.
\end{chapteropening}

\section[\texorpdfstring{\hi{घोलना} (gholnā)}{घोलना (gholnā)}]{\hi{घोलना} (gholnā) — to dissolve; to stir in}

\label{lesson:HI-C489-gholna}



\begin{quote}
Before the new one: say the Hindi for to have (something) looked at, then the Hindi for to have something done.
\end{quote}

\subsection*{You'll want to know: \hi{घोलना}}
\textbf{\hi{घोलना}} — \emph{gholnā} — \textquotedblleft{}to dissolve; to stir in\textquotedblright{}.

\begin{quote}
\hi{मीरा} \hi{ने} \hi{पानी} \hi{में} \hi{दवा} \hi{घोली।}
\end{quote}

\begin{quote}
\emph{Mīrā ne pānī meṁ davā gholī.}
\end{quote}

\begin{quote}
Mira dissolved the medicine in water.
\end{quote}

Past, with \hi{ने} (he or she …): \hi{उसने} \hi{दवा} \hi{घोली}, \emph{usne davā gholī}.

\begin{cousinweb}
\textbf{\hi{घोलना}} (\emph{gholnā}) is a tadbhava: CDIAL lists Hindi \emph{gholnā}, \textquotedblleft{}to mix, dissolve, melt\textquotedblright{}, under Sanskrit \emph{ghōlayati} (CDIAL 4526).
\end{cousinweb}

\begin{grammarlens}[title={What you've built}]
One more word for describing what you lived and felt.
\end{grammarlens}

\subsection*{Your turn}
\begin{itemize}
\raggedright
  \item \emph{Say it:} \emph{gholnā}
  \item \emph{Say it:} \emph{gholnā}, once more
  \item \emph{Say it:} \emph{Mīrā ne pānī meṁ davā gholī.} — Mira dissolved the medicine in water.
  \item \emph{Recall:} say the Hindi for to have (something) looked at, then the Hindi for to have something done, then say \emph{gholnā} again
\end{itemize}

\begin{tcolorbox}[breakable,colback=teal!4,colframe=teal!35!black,title={Before you move on}]
What does \hi{घोलना} mean? (\textquotedblleft{}To dissolve; to stir in\textquotedblright{}.) Say it once more.
\end{tcolorbox}
\section[\texorpdfstring{\hi{चटाना} (caṭānā)}{चटाना (caṭānā)}]{\hi{चटाना} (caṭānā) — to give (a baby) something to lick}

\label{lesson:HI-C489-catana}



\begin{quote}
Before the new one: say the Hindi for to have something done, then the Hindi for to dissolve.
\end{quote}

\subsection*{You'll want to know: \hi{चटाना}}
\textbf{\hi{चटाना}} — \emph{caṭānā} — \textquotedblleft{}to give (a baby) something to lick\textquotedblright{}.

\begin{quote}
\hi{माँ} \hi{ने} \hi{बच्चे} \hi{को} \hi{थोड़ा} \hi{शहद} \hi{चटाया।}
\end{quote}

\begin{quote}
\emph{māṁ ne bacce ko thoṛā śahad caṭāyā.}
\end{quote}

\begin{quote}
Mother gave the baby a little honey to lick.
\end{quote}

Past, with \hi{ने} (he or she …): \hi{उसने} \hi{बच्चे} \hi{को} \hi{दवा} \hi{चटाई}, \emph{usne bacce ko davā caṭāī}.

\begin{cousinweb}
\textbf{\hi{चटाना}} (\emph{caṭānā}) is the causative partner of \emph{cāṭnā}, \textquotedblleft{}to lick\textquotedblright{}, which CDIAL lists under \emph{*caṭṭ}, a form it reconstructs (CDIAL 4573). Medicine for small children is often given this way, mixed with honey.
\end{cousinweb}

\begin{grammarlens}[title={What you've built}]
One more word for describing what you lived and felt.
\end{grammarlens}

\subsection*{Your turn}
\begin{itemize}
\raggedright
  \item \emph{Say it:} \emph{caṭānā}
  \item \emph{Say it:} \emph{caṭānā}, once more
  \item \emph{Say it:} \emph{māṁ ne bacce ko thoṛā śahad caṭāyā.} — Mother gave the baby a little honey to lick.
  \item \emph{Recall:} say the Hindi for to have something done, then the Hindi for to dissolve, then say \emph{caṭānā} again
\end{itemize}

\begin{tcolorbox}[breakable,colback=teal!4,colframe=teal!35!black,title={Before you move on}]
What does \hi{चटाना} mean? (\textquotedblleft{}To give (a baby) something to lick\textquotedblright{}.) Say it once more.
\end{tcolorbox}
\section[\texorpdfstring{\hi{ख़ुराक} (khurāk)}{ख़ुराक (khurāk)}]{\hi{ख़ुराक} (khurāk) — a dose}

\label{lesson:HI-C489-khurak}



\begin{quote}
Before the new one: say the Hindi for to dissolve, then the Hindi for to give (a baby) something to lick.
\end{quote}

\subsection*{You'll want to know: \hi{ख़ुराक}}
\textbf{\hi{ख़ुराक}} — \emph{khurāk} — \textquotedblleft{}a dose\textquotedblright{}.

\begin{quote}
\hi{दवा} \hi{की} \hi{पहली} \hi{ख़ुराक} \hi{सुबह} \hi{लेनी} \hi{है।}
\end{quote}

\begin{quote}
\emph{davā kī pahlī khurāk subah lenī hai.}
\end{quote}

\begin{quote}
The first dose of the medicine is to be taken in the morning.
\end{quote}

Feminine.

\begin{cousinweb}
\textbf{\hi{ख़ुराक}} (\emph{khurāk}) is usually traced to Persian \emph{khurāk}, \textquotedblleft{}food, diet\textquotedblright{}; Hindi uses it for a measured amount of food or medicine, a dose.
\end{cousinweb}

\begin{grammarlens}[title={What you've built}]
One more word for describing what you lived and felt.
\end{grammarlens}

\subsection*{Your turn}
\begin{itemize}
\raggedright
  \item \emph{Say it:} \emph{khurāk}
  \item \emph{Say it:} \emph{khurāk}, once more
  \item \emph{Say it:} \emph{davā kī pahlī khurāk subah lenī hai.} — The first dose of the medicine is to be taken in the morning.
  \item \emph{Recall:} say the Hindi for to dissolve, then the Hindi for to give (a baby) something to lick, then say \emph{khurāk} again
\end{itemize}

\begin{tcolorbox}[breakable,colback=teal!4,colframe=teal!35!black,title={Before you move on}]
What does \hi{ख़ुराक} mean? (\textquotedblleft{}A dose\textquotedblright{}.) Say it once more.
\end{tcolorbox}
\section[\texorpdfstring{\hi{मरहम} (marham)}{मरहम (marham)}]{\hi{मरहम} (marham) — an ointment}

\label{lesson:HI-C489-marham}



\begin{quote}
Before the new one: say the Hindi for to give (a baby) something to lick, then the Hindi for a dose.
\end{quote}

\subsection*{You'll want to know: \hi{मरहम}}
\textbf{\hi{मरहम}} — \emph{marham} — \textquotedblleft{}an ointment\textquotedblright{}.

\begin{quote}
\hi{मरहम} \hi{से} \hi{घाव} \hi{को} \hi{जल्दी} \hi{आराम} \hi{मिला।}
\end{quote}

\begin{quote}
\emph{marham se ghāv ko jaldī ārām milā.}
\end{quote}

\begin{quote}
The ointment soon brought relief to the wound.
\end{quote}

Masculine.

\begin{cousinweb}
\textbf{\hi{मरहम}} (\emph{marham}) is usually traced to Arabic \emph{marham}, \textquotedblleft{}ointment\textquotedblright{}, which reached Hindi through Persian.
\end{cousinweb}

\begin{grammarlens}[title={What you've built}]
One more word for describing what you lived and felt.
\end{grammarlens}

\subsection*{Your turn}
\begin{itemize}
\raggedright
  \item \emph{Say it:} \emph{marham}
  \item \emph{Say it:} \emph{marham}, once more
  \item \emph{Say it:} \emph{marham se ghāv ko jaldī ārām milā.} — The ointment soon brought relief to the wound.
  \item \emph{Recall:} say the Hindi for to give (a baby) something to lick, then the Hindi for a dose, then say \emph{marham} again
\end{itemize}

\begin{tcolorbox}[breakable,colback=teal!4,colframe=teal!35!black,title={Before you move on}]
What does \hi{मरहम} mean? (\textquotedblleft{}An ointment\textquotedblright{}.) Say it once more.
\end{tcolorbox}
\section[\texorpdfstring{\hi{काढ़ा} (kāṛhā)}{काढ़ा (kāṛhā)}]{\hi{काढ़ा} (kāṛhā) — a herbal brew; a decoction}

\label{lesson:HI-C489-karha}



\begin{quote}
Before the new one: say the Hindi for a dose, then the Hindi for an ointment.
\end{quote}

\subsection*{You'll want to know: \hi{काढ़ा}}
\textbf{\hi{काढ़ा}} — \emph{kāṛhā} — \textquotedblleft{}a herbal brew; a decoction\textquotedblright{}.

\begin{quote}
\hi{ज़ुकाम} \hi{में} \hi{दादी} \hi{अदरक} \hi{का} \hi{काढ़ा} \hi{बनाती} \hi{हैं।}
\end{quote}

\begin{quote}
\emph{zukām meṁ dādī adrak kā kāṛhā banātī haiṁ.}
\end{quote}

\begin{quote}
When someone has a cold, my grandmother makes a ginger brew.
\end{quote}

Masculine.

\begin{cousinweb}
\textbf{\hi{काढ़ा}} (\emph{kāṛhā}) is probably the noun of \emph{kāṛhnā}, \textquotedblleft{}to pull out\textquotedblright{}, which CDIAL lists under \emph{*kaḍḍhati}, a form it reconstructs (CDIAL 2660): what boiling draws out of herbs. The link of senses is not given by CDIAL.
\end{cousinweb}

\begin{grammarlens}[title={What you've built}]
One more word for describing what you lived and felt.
\end{grammarlens}

\subsection*{Your turn}
\begin{itemize}
\raggedright
  \item \emph{Say it:} \emph{kāṛhā}
  \item \emph{Say it:} \emph{kāṛhā}, once more
  \item \emph{Say it:} \emph{zukām meṁ dādī adrak kā kāṛhā banātī haiṁ.} — When someone has a cold, my grandmother makes a ginger brew.
  \item \emph{Recall:} say the Hindi for a dose, then the Hindi for an ointment, then say \emph{kāṛhā} again
\end{itemize}

\begin{tcolorbox}[breakable,colback=teal!4,colframe=teal!35!black,title={Before you move on}]
What does \hi{काढ़ा} mean? (\textquotedblleft{}A herbal brew; a decoction\textquotedblright{}.) Say it once more.
\end{tcolorbox}
